\chapter{Episiotomy and Perineal Repair}

\section*{Introduction \& Clinical Indications}
An episiotomy is a surgical incision made in the perineum, the area between the vaginal opening and the anus, performed during the second stage of labor to facilitate childbirth. This procedure aims to enlarge the vaginal outlet, particularly in cases where there is a risk of severe perineal tearing or when rapid delivery is necessary due to fetal distress. The incision is typically made to prevent more complex lacerations that could lead to significant morbidity for the mother. Following the delivery, perineal repair involves the surgical closure of any perineal trauma, whether resulting from an episiotomy or spontaneous laceration. The primary objective of this repair is to restore the anatomical integrity of the perineum and pelvic floor, achieve hemostasis, and minimize the risk of long-term complications such as chronic pain, dyspareunia, and pelvic floor dysfunction, including urinary or fecal incontinence.

\begin{itemize}
  \item Perineotomy (for the incision)
  \item Obstetric Perineal Laceration Repair
  \item Vaginal Outlet Incision and Repair
  \item Perineorrhaphy (specifically referring to the surgical repair of the perineum)
  \item Obstetric Anal Sphincter Injury (OASIS) Repair (for third- and fourth-degree tears)
  \item Perineal Trauma Repair
\end{itemize}

The fundamental principle behind an episiotomy is the creation of a controlled, clean surgical incision in the perineum, which is theoretically easier to repair and less prone to complications than an uncontrolled, jagged, and potentially more extensive spontaneous tear. By strategically enlarging the birth canal opening, an episiotomy aims to expedite delivery in situations where fetal well-being is compromised or maternal expulsive efforts are insufficient, thereby reducing the duration of the second stage of labor and potential fetal hypoxia. The incision typically involves the skin, subcutaneous tissue, perineal body, and vaginal mucosa, with the specific direction (midline or mediolateral) chosen based on clinical assessment and risk factors.

Following delivery, the paramount principle of perineal repair is the meticulous, layered reconstruction of all injured tissues to restore anatomical and functional integrity. This involves the precise identification and reapproximation of the vaginal mucosa, the underlying perineal muscles (such as the bulbocavernosus, superficial transverse perineal muscles, and potentially the levator ani complex), and finally the subcutaneous tissue and skin. For more severe lacerations involving the anal sphincter (third-degree) or rectal mucosa (fourth-degree), the repair extends to these critical structures, aiming to restore their anatomical and functional integrity to prevent devastating long-term complications like fecal incontinence. The primary technical goals of repair are to achieve complete hemostasis, eliminate dead space, ensure optimal functional outcomes, minimize postoperative pain, and promote rapid, uncomplicated wound healing.

The practice of perineal incision during childbirth has a long and dynamic history. Early descriptions of perineal incisions to facilitate delivery can be traced back to the 18th century, notably by the Irish obstetrician Sir Fielding Ould in 1742. However, it was Sir James Young Simpson, a prominent Scottish obstetrician, who popularized the procedure in the mid-19th century, advocating for its use in cases of difficult labor. By the early 20th century, episiotomy became a widely adopted, often routine, practice in many parts of the world, particularly in North America. This widespread adoption was driven by the prevailing belief that it prevented severe perineal tearing, protected the pelvic floor from future prolapse, and reduced the risk of fetal intracranial hemorrhage.

However, by the latter half of the 20th century, a growing body of evidence began to challenge the routine application of episiotomy. Landmark studies demonstrated that routine episiotomy did not consistently prevent severe tears, nor did it offer significant protection against pelvic floor dysfunction or future prolapse. In fact, it was often associated with increased maternal morbidity, including greater pain, higher rates of infection, and paradoxically, an increased risk of third- and fourth-degree lacerations compared to spontaneous tearing. This paradigm shift led to major professional organizations, such as the American College of Obstetricians and Gynecologists (ACOG) and the Royal College of Obstetricians and Gynaecologists (RCOG), advocating for a restrictive use of episiotomy, reserving it for specific clinical indications rather than routine application. Concurrently, the understanding of perineal anatomy, particularly the anal sphincter complex, and the techniques for repairing obstetric anal sphincter injuries (OASIS) have advanced significantly, leading to improved surgical outcomes and reduced long-term morbidity.

An episiotomy is performed only when there is a clear, evidence-based clinical indication to facilitate a safe and timely vaginal delivery, while perineal repair is indicated for any degree of perineal trauma sustained during childbirth.

\begin{itemize}
  \item To expedite delivery in cases of non-reassuring fetal heart rate patterns or other signs of acute fetal distress, where rapid delivery is paramount.
  \item To facilitate operative vaginal delivery, such as with the use of forceps or vacuum extraction, which often require additional space for instrument placement and maneuverability.
  \item To prevent severe, irregular perineal lacerations when the perineum is judged to be unyielding, rigid, or excessively scarred, and tearing is imminent despite adequate perineal support.
  \item To aid in specific maneuvers required for the management of shoulder dystocia, where an enlarged vaginal opening may facilitate access and rotation of the fetal shoulders.
  \item To accommodate a suspected macrosomic infant, where the fetal head or body size is disproportionately large for the maternal pelvis, increasing the risk of severe tearing.
  \item To assist in delivery for a prolonged second stage of labor with documented maternal exhaustion and minimal progress, especially when other interventions have failed.
  \item To address a persistent occiput posterior position, which can lead to a wider diameter presenting to the perineum and a higher risk of complex perineal trauma.
  \item To repair any degree of perineal laceration (first, second, third, or fourth degree) or an episiotomy after delivery, which is a universal indication for anatomical restoration and prevention of complications.
\end{itemize}

An episiotomy is considered a **secondary obstetric intervention**, meaning it is not a routine or primary step in vaginal delivery but rather a procedure performed only when specific, well-defined indications arise during the labor process. It is an elective surgical intervention undertaken to facilitate delivery or prevent anticipated complications, and its use has become highly restrictive based on current evidence. It serves as an adjunct to the natural birthing process, employed only when the benefits of rapid delivery or prevention of severe, uncontrolled tearing outweigh the risks associated with the incision itself.

In contrast, the repair of perineal lacerations, whether spontaneous or resulting from an episiotomy, is a **primary and essential procedure** performed immediately after the delivery of the placenta. This repair is fundamental to postpartum care for nearly all women who experience perineal trauma during vaginal birth. Its primary role is to restore anatomical integrity, achieve hemostasis, prevent infection, and minimize long-term morbidity such as chronic pain, dyspareunia, and pelvic floor dysfunction. Therefore, while episiotomy is a secondary, indicated intervention, perineal repair is a fundamental and often unavoidable component of postpartum management following vaginal delivery.

\section*{Relevant Surgical Anatomy}
The perineum is a diamond-shaped region located inferior to the pelvic diaphragm, divided into an anterior urogenital triangle and a posterior anal triangle. Central to obstetric trauma and repair is the **perineal body**, a fibromuscular mass situated at the junction of these two triangles, between the posterior vaginal wall and the anus. This structure serves as a crucial attachment point for several muscles, including:

\begin{itemize}
  \item \textbf{Bulbocavernosus muscles:} Encircle the vaginal orifice and clitoris, contributing to vaginal tone.
  \item \textbf{Superficial transverse perineal muscles:} Extend from the ischial tuberosities to the perineal body.
  \item \textbf{Deep transverse perineal muscles:} Lie superior to the superficial muscles, forming part of the urogenital diaphragm.
  \item \textbf{External anal sphincter (EAS):} A voluntary striated muscle surrounding the anal canal, with superficial and deep portions that insert into the perineal body.
  \item \textbf{Levator ani muscles:} While not directly part of the perineum, the puborectalis portion of the levator ani forms a sling around the rectum, contributing significantly to fecal continence and pelvic floor support, and can be involved in deeper lacerations.
\end{itemize}

The **internal anal sphincter (IAS)**, an involuntary smooth muscle, lies deep to the EAS and is crucial for resting anal tone. The primary nerve supply to the perineum is the **pudendal nerve** (S2-S4), which provides sensory innervation to the external genitalia and perineum, and motor innervation to the external anal sphincter and perineal muscles. The main arterial supply is from branches of the **internal pudendal artery**, particularly the perineal artery, which provides rich vascularization. Venous drainage parallels the arterial supply, primarily via the internal pudendal veins. Lymphatic drainage from the superficial perineum typically drains to the superficial inguinal lymph nodes, while deeper structures may drain to internal iliac nodes. Surgical landmarks include the posterior fourchette, the midline raphe, and the ischial tuberosities, which guide the direction of episiotomy and subsequent repair.

\section*{Preoperative Workup \& Preparation}
Preparation for an episiotomy and perineal repair primarily occurs during the active phase of labor and immediately postpartum, as these are acute intrapartum events. Prior to delivery, during the prenatal period, the patient's medical history is reviewed for any conditions that might affect healing or increase bleeding risk, such as known coagulopathies, severe anemia, or inflammatory bowel disease. Informed consent for potential episiotomy and repair of any perineal lacerations is typically obtained during the prenatal period or upon admission to labor and delivery, as part of the general consent for childbirth.

Specific preparations at the time of delivery include:

\begin{itemize}
  \item \textbf{Anesthesia Assessment:} Ensuring adequate pain control is paramount. This involves assessing the efficacy of existing regional anesthesia (epidural/spinal) or preparing for local anesthetic infiltration of the perineum with agents like lidocaine.
  \item \textbf{Sterile Field Preparation:} The perineal area is meticulously cleansed with an antiseptic solution (e.g., povidone-iodine or chlorhexidine) and sterile drapes are applied to minimize the risk of wound infection.
  \item \textbf{Optimal Visualization:} Ensuring excellent lighting and exposure of the perineum is critical to accurately assess the extent of any trauma and perform a meticulous, layered repair.
  \item \textbf{Equipment Readiness:} Having appropriate sterile surgical instruments readily available, including sharp scissors, needle holders, tissue forceps, absorbable suture material of various sizes (e.g., polyglactin 2-0, 3-0, 4-0), sponges, and local anesthetic.
  \item \textbf{Comprehensive Assessment:} A thorough visual and digital examination of the perineum, including a rectovaginal examination, is performed immediately after delivery to accurately classify the degree of laceration and identify all injured structures, especially the anal sphincter and rectal mucosa.
  \item \textbf{Prophylactic Antibiotics:} For third- and fourth-degree perineal lacerations (Obstetric Anal Sphincter Injuries - OASIS), a single dose of broad-spectrum prophylactic antibiotics (e.g., a second-generation cephalosporin like cefazolin) is often administered intravenously to significantly reduce the risk of wound infection.
\end{itemize}

No specific NPO (nil per os) status or extensive cardiac workup is typically required beyond routine labor management, as these procedures are performed acutely.

\textbf{Absolute Contraindications for Episiotomy:}

\begin{itemize}
  \item \textbf{Absence of clear clinical indication:} Routine or prophylactic episiotomy is strongly discouraged and considered contraindicated due to overwhelming evidence of increased maternal morbidity without demonstrable benefit.
  \item \textbf{Active perineal infection:} Performing an episiotomy through an infected area (e.g., herpes simplex lesion, abscess) could worsen the infection, impair healing, and lead to systemic spread.
\end{itemize}

\textbf{Relative Contraindications and Precautions for Episiotomy:}

\begin{itemize}
  \item \textbf{Coagulopathy or bleeding disorder:} Increases the risk of excessive hemorrhage, hematoma formation, and difficulty achieving hemostasis during repair.
  \item \textbf{Inflammatory bowel disease (e.g., Crohn's disease) involving the perineum:} Significantly increases the risk of poor wound healing, fistula formation, and severe infection.
  \item \textbf{Previous severe perineal trauma with poor healing or rectovaginal fistula:} May warrant careful consideration of alternative delivery methods or a highly individualized approach to minimize recurrence.
  \item \textbf{Midline episiotomy in cases of short perineum or high risk of third- or fourth-degree tear:} A mediolateral episiotomy may be preferred to reduce the risk of anal sphincter involvement, despite its own drawbacks.
  \item \textbf{Preterm delivery:} The perineum of a preterm infant may be more fragile and prone to tearing, requiring careful assessment.
\end{itemize}

\textbf{Precautions for Perineal Repair:}

\begin{itemize}
  \item \textbf{Meticulous anatomical identification:} Precise identification of all tissue layers, especially the torn ends of the anal sphincter and the integrity of the rectal mucosa, is paramount to ensure proper repair and prevent long-term complications.
  \item \textbf{Strict aseptic technique:} Adherence to sterile technique is essential throughout the repair to minimize the risk of wound infection, particularly for deeper tears.
  \item \textbf{Adequate anesthesia:} Ensuring the patient is comfortable and pain-free throughout the repair is crucial for cooperation, a thorough procedure, and optimal outcomes.
  \item \textbf{Avoidance of excessive tension:} Sutures should approximate tissues without undue tension, which can lead to ischemia, increased pain, wound breakdown, and poor healing.
  \item \textbf{Post-repair rectal examination:} A final digital rectal examination is mandatory after repair of any third- or fourth-degree laceration to confirm the integrity of the rectal mucosa and ensure no sutures have inadvertently penetrated the rectal lumen.
\end{itemize}

\section*{Surgical Technique \& Procedure}
The procedure for episiotomy and perineal repair begins with ensuring adequate anesthesia. If the patient does not already have regional anesthesia (epidural or spinal block), local infiltration of the perineum with a local anesthetic such as 1\% lidocaine (often with epinephrine for vasoconstriction and prolonged effect) is performed. The patient is positioned in the dorsal lithotomy position, with legs supported in stirrups, allowing for optimal visualization and access to the perineum. The perineal area is then cleansed with an antiseptic solution and draped to establish a sterile field.

For an episiotomy, the incision is typically made with surgical scissors at the peak of a contraction, when the fetal head is maximally distending the perineum. The two main types are:

\begin{itemize}
  \item \textbf{Midline episiotomy:} A straight incision made from the posterior fourchette directly towards the anus. This type is generally easier to repair, less painful postoperatively, and heals well, but carries a higher risk of extending into a third- or fourth-degree laceration involving the anal sphincter or rectum.
  \item \textbf{Mediolateral episiotomy:} An incision made from the posterior fourchette at a 45- to 60-degree angle, typically to the right or left. This type is more difficult to repair, may be associated with more postoperative pain, and can result in greater blood loss, but it has a significantly lower risk of extending to involve the anal sphincter.
\end{itemize}

After the delivery of the infant and placenta, the perineum is thoroughly inspected to assess the extent of any laceration or episiotomy. A rectovaginal examination is crucial to rule out occult third- or fourth-degree extensions. The repair is then performed in layers using absorbable sutures, which do not require removal.

The key operative steps for perineal repair are:

\begin{enumerate}
  \item \textbf{Anesthesia and Inspection:} Ensure adequate local or regional anesthesia. Thoroughly inspect the entire perineum, vagina, and rectum to identify the full extent and depth of the laceration, classifying it (first, second, third, or fourth degree).
  \item \textbf{Rectal Mucosa Repair (for Fourth-Degree Tears):} If the rectal mucosa is involved, it is repaired first with fine, interrupted sutures (e.g., 3-0 or 4-0 polyglactin) that do not penetrate the rectal lumen, ensuring the knots are tied within the rectal wall.
  \item \textbf{Anal Sphincter Repair (for Third- and Fourth-Degree Tears):} The torn ends of the external anal sphincter (and internal anal sphincter if identified) are identified. The EAS is typically repaired using either an end-to-end (approximation) technique or an overlapping technique, with interrupted sutures (e.g., 2-0 or 3-0 polyglactin).
  \item \textbf{Vaginal Mucosa Repair:} The vaginal mucosa and submucosa are repaired next, starting approximately 1 cm above the apex of the laceration, using a continuous locking or non-locking suture (e.g., 2-0 polyglactin) to achieve hemostasis and reapproximate the vaginal wall.
  \item \textbf{Perineal Muscle Repair:} The severed or torn perineal muscles (e.g., bulbocavernosus, superficial transverse perineal muscles, deep transverse perineal muscles, and potentially levator ani) are identified and reapproximated with interrupted or figure-of-eight sutures (e.g., 2-0 polyglactin), restoring the perineal body.
  \item \textbf{Subcutaneous Tissue and Skin Closure:} The subcutaneous tissue is approximated to eliminate dead space. The final skin layer is closed with a subcuticular suture (e.g., 3-0 polyglactin) for optimal cosmetic results and to minimize skin tension, or with interrupted sutures.
  \item \textbf{Final Rectal Examination:} A mandatory digital rectal examination is performed after the completion of any third- or fourth-degree repair to confirm the integrity of the rectal mucosa and ensure no sutures have inadvertently penetrated the rectal lumen.
\end{enumerate}

Hemostasis is ensured throughout the repair, and the final repair is checked for proper anatomical alignment and absence of tension. No drains or implants are typically used in routine perineal repair.

\section*{Complications \& Risks}
While generally safe, episiotomy and perineal repair carry several potential risks, ranging from minor discomfort to severe long-term complications.

\begin{itemize}
  \item \textbf{General Surgical Complications:}
  \begin{itemize}
    \item \textbf{Bleeding and hematoma formation:} Accumulation of blood under the skin or within the perineal tissues, causing significant pain, swelling, and potentially requiring surgical evacuation.
    \item \textbf{Infection:} Perineal cellulitis, abscess formation, or wound breakdown, which can delay healing and increase morbidity.
    \item \textbf{Anesthesia complications:} Local anesthetic toxicity, allergic reactions, or side effects from regional anesthesia (e.g., headache, nerve injury).
  \end{itemize}
  \item \textbf{Procedure-Specific Risks:}
  \begin{itemize}
    \item \textbf{Increased perineal pain:} Both acute and chronic pain, potentially lasting for weeks or months postpartum, significantly impacting daily activities and quality of life.
    \item \textbf{Dyspareunia:} Painful sexual intercourse, which can persist long-term and affect intimate relationships.
    \item \textbf{Wound dehiscence:} Breakdown or separation of the repaired wound edges, often due to infection, tension, or poor healing, requiring re-suturing or secondary healing.
    \item \textbf{Extension of episiotomy:} An episiotomy, particularly a midline one, can extend spontaneously into a more severe third- or fourth-degree laceration involving the anal sphincter or rectum, despite the intention to prevent this.
    \item \textbf{Anal incontinence:} Inability to control flatus, liquid stool, or solid stool, especially after third- or fourth-degree tears, due to anal sphincter damage or inadequate repair. This can significantly impact quality of life.
    \item \textbf{Rectovaginal fistula:} A rare but severe complication of fourth-degree tears, where an abnormal connection forms between the rectum and vagina, leading to continuous fecal leakage into the vagina.
    \item \textbf{Scarring and keloid formation:} Excessive or hypertrophic scar tissue that can be painful, itchy, or cosmetically undesirable.
    \item \textbf{Nerve damage:} Rare injury to the pudendal nerve or its branches, leading to sensory deficits (numbness) or motor deficits (weakness) in the perineal area, potentially affecting continence or sexual function.
    \item \textbf{Incomplete healing or persistent discomfort:} Despite technically adequate repair, some women experience ongoing issues such as perineal heaviness, pressure, or discomfort.
  \end{itemize}
\end{itemize}

\section*{Postoperative Care \& Recovery}
Immediately after the episiotomy and perineal repair, the patient is monitored in the recovery area (PACU or labor and delivery suite) for vital signs, uterine tone, and signs of excessive bleeding or hematoma formation. Pain assessment and initial pain management strategies are implemented, typically involving oral analgesics. The perineal area is often covered with an ice pack for the first 12-24 hours to reduce swelling and discomfort. Early ambulation is encouraged to promote circulation and prevent complications like deep vein thrombosis.

Over the first 24-48 hours, pain management remains a priority, typically involving a combination of oral analgesics (e.g., NSAIDs like ibuprofen, acetaminophen, sometimes short-term opioids), topical anesthetic sprays or creams, and continued use of ice packs. Perineal hygiene is crucial; patients are instructed on using a peri bottle with warm water after urination and defecation, and frequent pad changes. Sitz baths (warm water soaks) can provide comfort, reduce swelling, and promote healing. Stool softeners are routinely prescribed to prevent straining during bowel movements, which could put undue stress on the repair, especially for third- and fourth-degree tears. Over the first 1-2 weeks, the absorbable sutures gradually dissolve, and the wound begins to heal. Pain typically subsides significantly, though some discomfort, especially with sitting, walking, or bowel movements, is common. Patients are advised to avoid heavy lifting, strenuous activities, and tampon use. Long-term recovery involves the complete resolution of pain, restoration of normal bowel and bladder function, and eventual return to comfortable sexual activity, usually after the 6-week postpartum check-up, once healing is complete and discomfort has resolved. Full functional recovery, particularly for severe tears, can take several months.

During an episiotomy and subsequent perineal repair, the surgeon's primary "findings" are the intraoperative observations of the extent and depth of the perineal trauma. This includes the classification of the laceration (first, second, third, or fourth degree) and the specific anatomical structures involved. For a first-degree laceration, only the vaginal mucosa and perineal skin are affected. A second-degree tear extends into the perineal body muscles but spares the anal sphincter. A third-degree laceration involves the anal sphincter complex (external and/or internal anal sphincter), while a fourth-degree tear extends through the anal sphincter into the rectal mucosa. The surgeon meticulously documents these findings, including the type of episiotomy (midline or mediolateral), the degree of any extension, the specific muscles identified and repaired, and the suture materials used.

Unlike many other surgical procedures, a formal "pathology report" from a laboratory is typically not generated for an episiotomy or perineal laceration repair, as no tissue is usually sent for histological examination unless there is an unusual finding such as a suspected tumor or an atypical infection. Therefore, the "pathology" in this context refers to the macroscopic anatomical disruption observed and documented by the operating surgeon. The detailed surgical note serves as the primary record of the findings and the repair performed, which is crucial for subsequent care and for understanding long-term outcomes.

A normal and successful postoperative course following an episiotomy and perineal repair is characterized by a progressive and predictable resolution of symptoms and restoration of function. Initially, patients can expect mild to moderate perineal pain and discomfort, which should gradually decrease over the first few days to weeks. The wound site will typically show signs of normal healing, such as mild swelling and redness, without excessive warmth, purulent discharge, or separation of the wound edges. Absorbable sutures will dissolve naturally, eliminating the need for removal. Patients should experience a gradual return to normal urinary and fecal continence, without leakage of urine, flatus, or stool. Bowel movements should become regular and comfortable with the aid of stool softeners. Most women are able to resume light daily activities within a few days and more strenuous activities, including exercise and sexual intercourse, typically after the 6-week postpartum check-up, once complete healing has occurred and discomfort has resolved. A cosmetically acceptable scar that is not excessively painful or hypertrophic is also an expected outcome. Ultimately, a normal course signifies minimal long-term morbidity and a full recovery from the physical trauma of childbirth.

Postoperative care after an episiotomy and perineal repair is crucial for monitoring healing, managing symptoms, and addressing any potential complications.

\begin{itemize}
  \item \textbf{Perineal Hygiene:} Patients are instructed to keep the perineal area clean and dry. This includes using a peri bottle with warm water to rinse after urination and defecation, patting dry gently, and changing sanitary pads frequently.
  \item \textbf{Pain Management:} Regular use of oral analgesics (NSAIDs, acetaminophen) and topical pain relief (sprays, creams) is encouraged. Ice packs for the first 24-48 hours and warm sitz baths thereafter can provide significant comfort.
  \item \textbf{Bowel Management:} Stool softeners are routinely prescribed to prevent constipation and straining, which can stress the perineal repair, especially for third- and fourth-degree tears.
  \item \textbf{Activity Restrictions:} Patients are advised to avoid heavy lifting, strenuous exercise, and prolonged standing or sitting for the first few weeks. Tampon use and sexual intercourse are typically deferred until after the 6-week postpartum check-up.
  \item \textbf{Postpartum Check-up:} A comprehensive postpartum visit is typically scheduled at 2 to 6 weeks after delivery. During this visit, the healthcare provider will visually inspect the perineal wound to assess healing, check for signs of infection or dehiscence, and evaluate for any persistent pain, dyspareunia, or incontinence.
  \item \textbf{Pelvic Floor Assessment:} For women who experienced third- or fourth-degree lacerations, a more thorough pelvic floor assessment may be performed, potentially including a digital rectal examination to assess anal sphincter tone and integrity.
  \item \textbf{Pelvic Floor Physical Therapy:} May be recommended for women experiencing persistent perineal pain, dyspareunia, urinary or fecal incontinence, or pelvic organ prolapse symptoms, to strengthen and rehabilitate the pelvic floor muscles.
  \item \textbf{Warning Signs:} Patients are educated on warning signs that require immediate medical attention, such as increasing perineal pain, fever, foul-smelling vaginal discharge, pus from the wound, wound separation, or new onset of urinary or fecal incontinence.
\end{itemize}

No specific imaging or laboratory tests are routinely performed unless complications are suspected.

\section*{Outcomes \& Prognosis}
The epidemiology of episiotomy and perineal lacerations has undergone substantial changes over recent decades, reflecting a global shift in obstetric practice. The incidence of episiotomy has dramatically declined in many developed countries, transitioning from a routine practice affecting over 60\% of vaginal deliveries in the 1970s and 80s to a restrictive use rate of typically less than 10-15\% today, in adherence to evidence-based guidelines. Despite this decline, perineal lacerations remain extremely common, with the vast majority of women (40-80\%) experiencing some degree of perineal trauma during vaginal birth. First- and second-degree lacerations are the most frequent, often requiring repair. More severe obstetric anal sphincter injuries (OASIS), which encompass third- and fourth-degree lacerations, occur in approximately 0.5\% to 5\% of vaginal deliveries, with rates varying based on population demographics, delivery practices, and reporting methods. Risk factors for OASIS include nulliparity, operative vaginal delivery (especially forceps), fetal macrosomia (birth weight >4000g), persistent occiput posterior position, and mediolateral episiotomy (though a well-timed mediolateral episiotomy can also prevent more severe tears in specific high-risk circumstances). These procedures are exclusively performed in women during childbirth. Morbidity associated with perineal trauma and repair includes acute pain, infection, wound dehiscence, and long-term complications such as dyspareunia and varying degrees of anal incontinence. Mortality directly attributable to episiotomy or perineal repair is exceedingly rare, usually linked to severe hemorrhage or sepsis from untreated complications.

The outcomes and prognosis following episiotomy and perineal repair are generally favorable, particularly for first- and second-degree lacerations. The success rate for primary repair of these less severe tears, in terms of anatomical restoration and symptom resolution, is very high, often exceeding 95\%. Most women experience good healing, resolution of pain within a few weeks, and a return to normal function, including comfortable sexual activity, typically by the 6-week postpartum mark.

For third- and fourth-degree lacerations (Obstetric Anal Sphincter Injuries, OASIS), the prognosis is more variable and depends significantly on the quality of the initial repair. While immediate and meticulous repair significantly improves outcomes, a substantial proportion of women may experience persistent symptoms. Success rates for restoring continence after OASIS repair range from 60-90\% for flatus and 40-70\% for solid stool, indicating that a significant minority may suffer from some degree of anal incontinence. Dyspareunia is also a common long-term complaint, affecting up to 20-30\% of women after severe tears. The risk of recurrence of a severe perineal tear in subsequent vaginal deliveries is approximately 4-8\%. Prognostic factors influencing outcomes include the initial severity of the tear, the quality and timing of the repair, the experience of the surgeon, the presence of postoperative complications like infection or wound dehiscence, and subsequent obstetric events. Despite potential complications, the vast majority of women recover well and lead normal lives, though some may require ongoing pelvic floor physical therapy or further surgical intervention for persistent issues.

\section*{References}
\begin{enumerate}
  \item MedlinePlus. Episiotomy. U.S. National Library of Medicine; 2024. Available from: \url{https://medlineplus.gov/episiotomy.html}
  \item Cunningham F, Leveno K, Bloom S, et al. Williams Obstetrics. 26th ed. New York, NY: McGraw-Hill Education; 2022.
  \item American College of Obstetricians and Gynecologists. ACOG Practice Bulletin No. 198: Prevention and Management of Obstetric Lacerations. Obstet Gynecol. 2018;132(3):e87-e102.
  \item Royal College of Obstetricians and Gynaecologists. Green-top Guideline No. 29: Management of Third- and Fourth-Degree Perineal Tears. London: RCOG; 2015.
\end{enumerate}

\clearpage
